\documentclass[../DoAn.tex]{subfiles}
\begin{document}

Chương này trình bày các công nghệ và nền tảng lý thuyết được sử dụng trong đồ án. Nội dung tập trung vào những khái niệm cần thiết để xây dựng trò chơi 2D, tổ chức mã nguồn, xử lý chuỗi tiếng Việt, sinh bảng ô chữ và lưu dữ liệu cục bộ. Các chi tiết thiết kế cụ thể của từng scene, script và luồng giao diện được trình bày ở Chương 4.

\section{Godot Engine}
Godot là game engine mã nguồn mở hỗ trợ phát triển trò chơi 2D và 3D, cung cấp hệ thống scene, node, signal, tài nguyên và công cụ dựng giao diện phục vụ quá trình phát triển trò chơi \cite{godotDocs}. Trong phạm vi đồ án, Godot được lựa chọn vì phù hợp với trò chơi 2D có nhiều thành phần giao diện, cần phản hồi nhanh với thao tác người chơi và không yêu cầu hệ thống vật lý phức tạp.

Một khái niệm quan trọng của Godot là scene. Scene có thể biểu diễn một màn hình hoàn chỉnh, một cụm giao diện hoặc một đối tượng tái sử dụng. Bên trong scene, các thành phần được tổ chức thành cây node. Node có thể là node giao diện như \code{Control}, \code{Button}, \code{Label}, \code{GridContainer}, hoặc node logic dùng để xử lý trạng thái và sự kiện. Cách tổ chức này phù hợp với trò chơi giải đố từ vựng dạng ô chữ vì bảng ô chữ, danh sách gợi ý, thanh thông tin và hộp thoại kết quả đều có thể được tách thành các cụm giao diện có trách nhiệm rõ ràng.

Godot cũng hỗ trợ signal để truyền sự kiện giữa các node. Signal giúp một node thông báo rằng một hành động đã xảy ra mà không cần biết trực tiếp toàn bộ logic xử lý phía sau. Ví dụ, một nút giao diện có thể phát signal khi được bấm; node điều phối nhận signal đó và quyết định chuyển màn, tạo màn chơi mới hoặc cập nhật trạng thái. Cơ chế này giúp giảm phụ thuộc trực tiếp giữa các thành phần giao diện và logic điều khiển.

\begin{figure}[H]
\centering
\resizebox{0.88\textwidth}{!}{\begin{tikzpicture}[
    font=\small,
    box/.style={draw, rounded corners, align=center, minimum width=3.0cm, minimum height=0.85cm, fill=white},
    scene/.style={draw, rounded corners, align=center, minimum width=10.5cm, minimum height=3.9cm, fill=gray!8},
    signal/.style={draw, rounded corners, align=center, minimum width=3.2cm, minimum height=0.85cm, fill=yellow!12},
    arrow/.style={-Stealth, thick}
]
\node[scene] (scene) at (5.8,0) {};
\node[font=\bfseries] at ($(scene.north)+(0,-0.35)$) {Scene};

\node[box] (control) at (2.2,0.75) {Control};
\node[box] (button) at (5.8,0.75) {Button};
\node[box] (label) at (9.4,0.75) {Label};
\node[box] (grid) at (5.8,-0.75) {GridContainer};
\node[box] (logic) at (9.4,-0.75) {Node logic};

\node[box, fill=blue!6] (player) at (-1.0,0.75) {Thao tác\\người chơi};
\node[signal] (sig) at (5.8,-2.65) {Signal\\sự kiện};
\node[box, fill=green!8] (handler) at (10.0,-2.65) {Logic xử lý};

\draw[arrow] (player.east) -- (control.west);
\draw[arrow] (control.east) -- (button.west);
\draw[arrow] (button.south) -- ++(0,-0.45) -- (sig.north);
\draw[arrow] (grid.south) -- ++(0,-0.45) -- (sig.north);
\draw[arrow] (sig.east) -- (handler.west);
\draw[arrow] (handler.north) -- (logic.south);
\draw[arrow, dashed] (logic.west) -- (grid.east);
\end{tikzpicture}
}
\caption{Mô hình scene, node và signal trong Godot}
\label{fig:godot-scene-node-signal}
\end{figure}

Trong triển khai của đồ án, các khái niệm scene, node và signal được áp dụng để tách phần menu, phần màn chơi và các hộp thoại. Tuy nhiên, Chương 3 chỉ trình bày vai trò công nghệ ở mức nền tảng; cấu trúc cụ thể của các scene và script được phân tích trong chương thiết kế và đánh giá hệ thống.

\section{GDScript}
GDScript là ngôn ngữ kịch bản được thiết kế cho Godot, có cú pháp gần với Python và tích hợp trực tiếp với hệ thống node của engine \cite{godotGDScript}. Việc sử dụng GDScript giúp mã nguồn có thể truy cập node, signal, tài nguyên và vòng đời scene một cách trực tiếp, không cần thêm lớp tích hợp trung gian như khi sử dụng một ngôn ngữ ngoài.

GDScript hỗ trợ khai báo lớp bằng \code{class\_name}, nhờ đó các thành phần logic có thể được đặt tên rõ ràng và sử dụng lại trong nhiều script. Với trò chơi giải đố từ vựng dạng ô chữ, cách tổ chức này phù hợp để tách các nhóm trách nhiệm như sinh bảng ô chữ, lưu dữ liệu, quản lý trạng thái màn chơi và chuẩn hóa chuỗi tiếng Việt. Khi các nhóm trách nhiệm được tách riêng, mã nguồn dễ đọc hơn và việc kiểm tra lỗi cũng thuận lợi hơn.

Godot cho phép kết hợp script gắn với node và lớp dữ liệu không nằm trong cây scene. Script gắn với node phù hợp cho các thành phần cần nhận sự kiện giao diện hoặc cập nhật hiển thị. Ngược lại, các lớp không phụ thuộc giao diện phù hợp cho xử lý dữ liệu, thuật toán và trạng thái. Sự phân tách này là cơ sở để Chương 4 trình bày kiến trúc ứng dụng theo hướng tách giao diện, trạng thái và xử lý nghiệp vụ.

\section{Xử lý chuỗi tiếng Việt}
Tiếng Việt sử dụng bảng chữ cái Latin mở rộng với dấu thanh và các biến thể nguyên âm. Trong Unicode, cùng một ký tự hiển thị có thể được biểu diễn bằng nhiều dạng mã hóa khác nhau, ví dụ ký tự dựng sẵn hoặc tổ hợp giữa ký tự cơ sở và dấu kết hợp. Unicode Standard Annex \#15 định nghĩa các dạng chuẩn hóa Unicode nhằm đưa chuỗi về dạng nhất quán trước khi xử lý \cite{unicodeNormalization}. Vì vậy, hệ thống cần có bước chuẩn hóa trước khi so khớp đáp án và sinh dữ liệu cho bảng ô chữ.

Đối với trò chơi giải đố từ vựng dạng ô chữ, một đáp án có thể được nhìn ở nhiều dạng khác nhau. Dạng hiển thị giữ nguyên dấu và khoảng trắng để người chơi đọc được nghĩa đầy đủ. Dạng dùng để xếp bảng cần bỏ khoảng trắng vì mỗi ô trên bảng ô chữ chỉ chứa một ký tự. Dạng dùng để so khớp đáp án cần chuẩn hóa chữ thường và xử lý khoảng trắng để hạn chế sai khác do cách nhập liệu. Việc tách các dạng biểu diễn này giúp hệ thống vừa giữ đúng tiếng Việt khi hiển thị, vừa xử lý được bài toán đặt từ trên bảng.

\begin{figure}[H]
\centering
\resizebox{0.90\textwidth}{!}{\begin{tikzpicture}[
    font=\small,
    flowstep/.style={draw, rounded corners, align=center, minimum width=3.3cm, minimum height=0.8cm, fill=white},
    data/.style={draw, rounded corners, align=center, minimum width=3.3cm, minimum height=0.8cm, fill=blue!6},
    arrow/.style={-Stealth, thick}
]
\node[data] (raw) at (0,0) {Đáp án gốc};
\node[flowstep] (lower) at (4.0,0) {Chuẩn hóa\\chữ thường};
\node[flowstep] (space) at (8.0,0) {Bỏ khoảng trắng\\khi xếp bảng};
\node[data] (puzzle) at (12.0,0) {Tạo\\\texttt{puzzle\_word}};

\node[flowstep] (split) at (12.0,-1.8) {Tách ký tự\\tiếng Việt};
\node[data] (parts) at (8.0,-1.8) {Lưu ký tự gốc,\\cơ sở, biến thể,\\thanh điệu};
\node[flowstep] (match) at (4.0,-1.8) {So khớp đáp án\\người chơi};
\node[data] (result) at (0,-1.8) {Kết quả\\đúng/sai};

\draw[arrow] (raw) -- (lower);
\draw[arrow] (lower) -- (space);
\draw[arrow] (space) -- (puzzle);
\draw[arrow] (puzzle) -- (split);
\draw[arrow] (split) -- (parts);
\draw[arrow] (parts) -- (match);
\draw[arrow] (match) -- (result);
\end{tikzpicture}
}
\caption{Quy trình xử lý chuỗi tiếng Việt trong hệ thống}
\label{fig:quy-trinh-xu-ly-tieng-viet}
\end{figure}

Trong quá trình xử lý, hệ thống cần lưu đồng thời thông tin ký tự gốc, ký tự cơ sở, biến thể nguyên âm và thanh điệu. Cách biểu diễn này giúp việc hiển thị vẫn giữ đúng dấu tiếng Việt, đồng thời tạo tiền đề cho những xử lý mở rộng như phân tích ký tự, hỗ trợ gợi ý theo chữ cái hoặc kiểm tra các biến thể nhập liệu. Phần triển khai cụ thể của bộ chuẩn hóa tiếng Việt được trình bày trong Chương 4.

\section{Sinh bảng ô chữ}
Sinh bảng ô chữ là bài toán đặt nhiều từ lên một lưới sao cho các từ có thể giao nhau tại những ký tự giống nhau, đồng thời không tạo ra các chuỗi phụ không hợp lệ. Với phạm vi đồ án, bài toán được tiếp cận theo hướng heuristic thay vì tìm tối ưu toàn cục. Cách tiếp cận này phù hợp vì yêu cầu chính của trò chơi là tạo được màn chơi hợp lệ trong thời gian chấp nhận được, bảng đủ gọn và các từ có liên kết với nhau.

Một quy trình sinh bảng ở mức khái niệm thường bắt đầu từ danh sách từ đã chuẩn hóa. Hệ thống chọn thứ tự thử, đặt từ đầu tiên gần trung tâm, sau đó tìm các vị trí giao nhau giữa từ mới và các từ đã đặt. Mỗi vị trí ứng viên cần được kiểm tra điều kiện hợp lệ: nằm trong phạm vi lưới, không mâu thuẫn ký tự đã có, không nối dài ngoài ý muốn và không tạo đoạn chữ phụ. Sau khi lọc các vị trí hợp lệ, hệ thống chấm điểm để ưu tiên vị trí có nhiều giao điểm, gần trung tâm và làm tăng diện tích bảng ít hơn.

\begin{figure}[H]
\centering
\resizebox{0.62\textwidth}{!}{\begin{tikzpicture}[
    font=\small,
    flowstep/.style={draw, rounded corners, align=center, minimum width=4.2cm, minimum height=0.72cm, fill=white},
    decision/.style={draw, diamond, aspect=2.2, align=center, inner sep=1pt, minimum width=3.1cm, fill=yellow!10},
    arrow/.style={-Stealth, thick},
    note/.style={font=\scriptsize}
]
\node[flowstep] (words) at (0,0) {Danh sách từ};
\node[flowstep, below=0.55cm of words] (sort) {Sắp xếp hoặc chọn thứ tự thử};
\node[flowstep, below=0.55cm of sort] (first) {Đặt từ đầu tiên gần trung tâm};
\node[flowstep, below=0.55cm of first] (cross) {Tìm vị trí giao nhau};
\node[decision, below=0.75cm of cross] (valid) {Vị trí\\hợp lệ?};
\node[flowstep, below=0.75cm of valid] (score) {Chấm điểm vị trí};
\node[flowstep, below=0.55cm of score] (place) {Cập nhật bảng};
\node[decision, below=0.75cm of place] (more) {Còn từ\\cần đặt?};
\node[flowstep, below=0.75cm of more] (output) {Trả về bảng và danh sách từ đã đặt};

\draw[arrow] (words) -- (sort);
\draw[arrow] (sort) -- (first);
\draw[arrow] (first) -- (cross);
\draw[arrow] (cross) -- (valid);
\draw[arrow] (valid) -- node[note, left] {Có} (score);
\draw[arrow] (score) -- (place);
\draw[arrow] (place) -- (more);
\draw[arrow] (more) -- node[note, left] {Không} (output);

\draw[arrow] (valid.east) -- ++(1.25,0) |- node[note, pos=0.18, right] {Không} (cross.east);
\draw[arrow] (more.east) -- ++(1.25,0) |- node[note, pos=0.18, right] {Có} (cross.east);
\end{tikzpicture}
}
\caption{Luồng sinh bảng ô chữ ở mức khái niệm}
\label{fig:luong-sinh-bang-o-chu}
\end{figure}

Hạn chế của heuristic là không bảo đảm tìm được phương án tốt nhất cho mọi tập từ. Một số danh sách từ có ít ký tự chung hoặc độ dài quá lệch có thể khó tạo thành bảng gọn. Tuy nhiên, cách tiếp cận này có ưu điểm là dễ triển khai, dễ kiểm soát điều kiện hợp lệ và có thể điều chỉnh bằng các tiêu chí chấm điểm. Chi tiết về cách áp dụng vào mã nguồn của đồ án được trình bày trong phần thiết kế và triển khai.

\section{Lưu dữ liệu cục bộ}
Ứng dụng cần lưu hồ sơ người chơi, tổng số sao, điểm cao nhất và tiến độ từng màn. Với chế độ chơi cục bộ, định dạng JSON là lựa chọn phù hợp vì dữ liệu của trò chơi chủ yếu gồm các cặp khóa - giá trị, danh sách và đối tượng lồng nhau. JSON là định dạng trao đổi dữ liệu dạng văn bản được chuẩn hóa trong RFC 8259 \cite{rfc8259}, dễ đọc, dễ ghi và được hỗ trợ trực tiếp trong nhiều môi trường phát triển.

Khi lưu dữ liệu cục bộ, hệ thống cần xử lý cả trường hợp tệp lưu chưa tồn tại, thiếu trường hoặc có cấu trúc không hợp lệ. Cách xử lý phù hợp là kiểm tra kiểu dữ liệu sau khi đọc, bổ sung giá trị mặc định cho trường bị thiếu và phục hồi cấu trúc hợp lệ nếu phát hiện dữ liệu lỗi. Điều này giúp ứng dụng không dừng đột ngột khi tệp lưu bị hỏng hoặc được tạo từ phiên bản cũ.

Lưu dữ liệu cục bộ có ưu điểm là đơn giản, không yêu cầu máy chủ và phù hợp với nguyên mẫu trò chơi chạy trên máy cá nhân. Hạn chế của cách làm này là dữ liệu khó đồng bộ giữa nhiều thiết bị. Trong phạm vi đồ án, lựa chọn lưu cục bộ là hợp lý vì mục tiêu chính là xây dựng trò chơi hoạt động ổn định, có tiến độ riêng cho từng hồ sơ người chơi trên cùng một máy.

\section{Công cụ phát triển và biên soạn báo cáo}
Quá trình phát triển sử dụng Godot 4.x cho phần ứng dụng trò chơi. Mã nguồn được chỉnh sửa bằng môi trường phát triển hỗ trợ đọc cấu trúc thư mục, tìm kiếm nhanh và quản lý nhiều file script. Báo cáo được biên soạn bằng LaTeX trên Overleaf, đồng thời có thể chỉnh sửa trực tiếp trong Visual Studio Code khi cần làm việc với toàn bộ project.

LaTeX phù hợp với báo cáo kỹ thuật vì hỗ trợ cấu trúc chương mục, đánh số hình bảng, trích dẫn tài liệu tham khảo và định dạng nhất quán. Overleaf giúp biên dịch tài liệu trực tuyến và thuận tiện khi cần chia sẻ bản báo cáo với giảng viên hướng dẫn. Việc chia báo cáo thành nhiều file theo từng chương cũng giúp quá trình chỉnh sửa rõ ràng hơn, tránh phải thao tác trên một file quá dài.

Các công nghệ và kỹ thuật trong chương này tạo thành nền tảng cho phần thiết kế và triển khai ứng dụng. Chương tiếp theo trình bày cách các thành phần đó được tổ chức thành kiến trúc cụ thể của trò chơi giải đố từ vựng dạng ô chữ.

\end{document}

